
% ============================================================
% CONTENIDOS DEL CURSO
%
% Esta tabla presenta la estructura conceptual reutilizable.
% No incluye semanas ni fechas: esos datos pertenecen al
% cronograma particular de cada edición.
% ============================================================

% Crea el encabezado numerado de la sección.
\section{Contenidos del curso}

Presenta aquí la organización conceptual del curso. Las semanas y
fechas específicas deben permanecer en el cronograma de cada
edición.


% ------------------------------------------------------------
% CONFIGURACIÓN LOCAL DE LA TABLA
% ------------------------------------------------------------

% \begingroup limita los ajustes siguientes únicamente a esta
% tabla. Al llegar a \endgroup, LaTeX recuperará los valores
% generales del documento.
\begingroup

% Reduce ligeramente el tamaño del texto de la tabla.
\small

% Aumenta la altura de las filas.
\renewcommand{\arraystretch}{1.25}

% Controla el espacio horizontal interno de las celdas.
\setlength{\tabcolsep}{7pt}


% ------------------------------------------------------------
% TABLA DE CONTENIDOS
% ------------------------------------------------------------

% tabularx crea una tabla que ocupa todo el ancho del texto.
%
% La primera columna:
% - tiene un ancho fijo del 22 %;
% - utiliza negrita;
% - alinea el contenido a la izquierda.
%
% La segunda columna:
% - usa X para ocupar el espacio restante;
% - alinea el contenido a la izquierda.
%
% @{} elimina el espacio automático de los bordes.
\begin{tabularx}{\textwidth}{@{}
    >{\bfseries\raggedright\arraybackslash}p{0.22\textwidth}
    >{\raggedright\arraybackslash}X
@{}}

% booktabs proporciona líneas horizontales tipográficas.
\toprule

% & separa las columnas.
% \\ termina la fila.
\textbf{Unidad o bloque} & \textbf{Contenidos} \\

% Separa el encabezado del cuerpo de la tabla.
\midrule

Unidad 1 &
Incorpora aquí los contenidos de la primera unidad. \\

% Añade espacio vertical sin dibujar otra línea.
\addlinespace[4pt]

Unidad 2 &
Incorpora aquí los contenidos de la segunda unidad. \\

% Cierra visualmente la tabla.
\bottomrule

\end{tabularx}

% Finaliza el grupo de ajustes locales.
\endgroup